\chapter{Liquidity and Funding Risk; Bank Treasury}\label{ch:rc:liquidity-and-funding-risk-bank-treasury}

On 8 March 2023 Silicon Valley Bank's parent announced that it had sold 21 billion dollars of
securities at an after-tax loss of 1.8 billion and would raise 2.25 billion of capital. The next
day depositors withdrew more than 40 billion dollars, and the bank's management expected more
than 100 billion to follow on 10 March, about 85\% of its deposits. The bank failed that day.
It had invested its deposits, 94\% of them uninsured, in long-dated agency mortgage securities
held to maturity, with a duration of 6.2 years; when rates rose, their value fell and its
depositors, a tightly connected network of technology firms and their investors, ran. Nothing in
the story involves a trading desk: it is the market risk of a bank's balance sheet and the
liquidity risk of its funding, which are managed by its treasury. This chapter measures both: the
regulatory liquidity ratios, the transfer prices that charge businesses for funding, the
interest rate risk of the \omterm{def:rc:regulatory-capital-for-trading-books:books}{banking book}, and the behaviour of deposits.

\section{Liquidity ratios}

\begin{definition}[Liquidity coverage ratio, high-quality liquid assets]\label{def:rc:liquidity-and-funding-risk-bank-treasury:lcr}
The \emph{liquidity coverage ratio}\index{liquidity coverage ratio} (LCR) is the stock of
\emph{high-quality liquid assets}\index{high-quality liquid assets} (HQLA) over the net cash outflows
of a 30-day stress scenario; it must be at least 100\%. HQLA are cash, central bank reserves and
Level~1 securities in full, Level~2A assets (such as agency mortgage securities) after a 15\% haircut,
and Level~2 assets are capped at 40\% of the stock; outflows apply run-off rates to each funding source,
from 5\% for insured stable retail deposits to 40\% for non-operational corporate deposits.
\end{definition}

\begin{definition}[Net stable funding ratio]\label{def:rc:liquidity-and-funding-risk-bank-treasury:nsfr}
The \emph{net stable funding ratio}\index{net stable funding ratio} (NSFR) is available stable
funding over required stable funding on a one-year horizon, weighted by how stable each liability is
and how illiquid each asset is; it must be at least 100\%.
\end{definition}

\begin{example}[A stylised bank]\label{ex:rc:liquidity-and-funding-risk-bank-treasury:bank}
An illustrative bank (USD billion) holds 15 of cash, 25 of five-year Treasuries, 90 of agency mortgage
securities (a level-pay schedule over twelve years, duration 6.3 at purchase) and 70 of floating-rate
loans, funded by 173 of deposits, 11 of term debt and 16 of equity. A year after rates rise from 1.0\%
to 4.5\%, its HQLA at market value are 61.0: 36.6 of Level~1 (cash and Treasuries) and 24.4 of Level~2A,
the MBS after the haircut and the cap. With 6\% of its deposits insured and stable (run-off 5\%) and
94\% uninsured (40\%), 30-day outflows are 65.6 and the LCR is 93\%, below the minimum.
\end{example}

\omcode{code/firm/alm/firm_alm.py}{122}{131}{The LCR: Level~1 and Level~2A assets with the haircut and the cap, against run-off-weighted outflows.}\label{lst:rc:liquidity-and-funding-risk-bank-treasury:lcr}

\section{Funds-transfer pricing}

\begin{definition}[Funds-transfer pricing]\label{def:rc:liquidity-and-funding-risk-bank-treasury:ftp}
\emph{Funds-transfer pricing}\index{funds-transfer pricing} (FTP) charges each business that uses
funding, and credits each that raises it, an internal rate: a base rate curve for the repricing
tenor plus a liquidity premium for the term of the funding, so that a desk's margin is measured
after the cost of the money it uses and the treasury bears (and hedges) the interest rate and
liquidity mismatch centrally.
\end{definition}

\begin{example}[An FTP curve]\label{ex:rc:liquidity-and-funding-risk-bank-treasury:ftp}
With a base rate of 4.5\% and an illustrative liquidity premium of $0.15\%\sqrt T$, a three-month loan is
charged 4.58\%, a five-year loan 4.84\% and a ten-year one 4.97\%; a deposit expected to stay five years
is credited 4.84\%. A business that funds long assets with short deposits sees the term premium in its
charge instead of in the treasury's losses.
\end{example}

\section{Asset--liability management and rate risk in the banking book}

\begin{definition}[Asset--liability management, IRRBB]\label{def:rc:liquidity-and-funding-risk-bank-treasury:alm}
\emph{Asset--liability management}\index{asset--liability management} (ALM) manages the mismatch
between a bank's assets and liabilities in rates, maturities and currencies. Its central market
risk is \emph{interest-rate risk in the banking book}\index{interest-rate risk in the banking book}
(IRRBB), measured in value and in earnings.
\end{definition}

\begin{definition}[Economic value of equity, NII sensitivity]\label{def:rc:liquidity-and-funding-risk-bank-treasury:eve}
The \emph{economic value of equity}\index{economic value of equity} (EVE) is the present value of the
\omterm{def:rc:regulatory-capital-for-trading-books:books}{banking book}'s asset cash flows less that of its liability cash flows (deposits slotted by
behavioural models). \emph{Net interest income sensitivity}\index{net interest income sensitivity} is
the change in the next year's interest income less interest expense under a rate scenario, on a
constant balance sheet.
\end{definition}

The supervisory standard prescribes six instantaneous shocks per currency: parallel up and down,
short rates up and down (largest at the short end, decaying as $e^{-t/4}$), a steepener and a
flattener, sized for the US dollar at 200, 300 and 150 basis points for parallel, short and long
moves (\cref{fig:rc:liquidity-and-funding-risk-bank-treasury:shocks}; the chart of \cref{fig:rc:liquidity-and-funding-risk-bank-treasury:eve} marks up moves $+$ and down moves $-$). A bank whose largest EVE loss
exceeds 15\% of its Tier~1 capital is an outlier.

\begin{dated}{2026-09}{The shock sizes}\label{dat:rc:liquidity-and-funding-risk-bank-treasury:shocks}
In July 2024 the Basel Committee recalibrated the shocks on data to December 2023, with local shock
factors, a 99.9th percentile and rounding to 25 basis points, for implementation by 1 January 2026. For the
US dollar the parallel and short shocks stay at 200 and 300 basis points and the long shock rises from 150
to 225. The chapter's figures use the 2016 sizes, in force in 2022.
\end{dated}

\begin{omfigure}
\begin{tikzpicture}
\begin{axis}[width=11.5cm,height=5.4cm,
  xlabel={tenor (years)},ylabel={shock (bp)},
  tick label style={font=\small},label style={font=\small},
  xmin=0,xmax=20,ymin=-320,ymax=320,grid=major,grid style={gray!25},
  legend columns=3,legend cell align=left,legend style={font=\footnotesize,at={(0.5,-0.25)},anchor=north,draw=none,fill=none,/tikz/every even column/.append style={column sep=8pt}},
  table/col sep=comma]
\addplot[omDef,very thick] table[x=t,y=pu]{figdata/rates-credit-risk/24-liquidity-and-funding-risk-bank-treasury/shocks.csv};
\addplot[omDef,very thick,dashed] table[x=t,y=pd]{figdata/rates-credit-risk/24-liquidity-and-funding-risk-bank-treasury/shocks.csv};
\addplot[omThm,thick] table[x=t,y=su]{figdata/rates-credit-risk/24-liquidity-and-funding-risk-bank-treasury/shocks.csv};
\addplot[omThm,thick,dashed] table[x=t,y=sd]{figdata/rates-credit-risk/24-liquidity-and-funding-risk-bank-treasury/shocks.csv};
\addplot[omMeth,thick,mark=*,mark size=1pt] table[x=t,y=st]{figdata/rates-credit-risk/24-liquidity-and-funding-risk-bank-treasury/shocks.csv};
\addplot[omExo,thick,mark=square*,mark size=1pt] table[x=t,y=fl]{figdata/rates-credit-risk/24-liquidity-and-funding-risk-bank-treasury/shocks.csv};
\legend{parallel up,parallel down,short up,short down,steepener,flattener}
\end{axis}
\end{tikzpicture}

\omcaption{The six supervisory interest rate shocks for the US dollar by tenor: parallel
$\pm200$ basis points, short-rate $\pm300$ decaying with tenor, and the two rotations built from the
short and long shocks. Data: Basel Committee IRRBB standard (2016); the chapter's
tutorial.}\label{fig:rc:liquidity-and-funding-risk-bank-treasury:shocks}
\end{omfigure}

\omcode{code/firm/alm/firm_alm.py}{19}{37}{The six supervisory interest rate shock scenarios.}\label{lst:rc:liquidity-and-funding-risk-bank-treasury:shocks}

\begin{example}[The stylised bank before the rise]\label{ex:rc:liquidity-and-funding-risk-bank-treasury:eve}
At a 1\% rate the bank's EVE is 22.9. The parallel-up shock cuts it by 6.95, 43\% of its Tier~1
capital of 16, nearly three times the outlier threshold; the steepener by 4.60 (6.79 with the 2024 sizes);
the parallel-down shock raises it by 8.64 (\cref{fig:rc:liquidity-and-funding-risk-bank-treasury:eve}). The risk was
measurable a year before it was realised.
\end{example}

\begin{omfigure}
\begin{tikzpicture}
\begin{axis}[width=11.5cm,height=5.2cm,ybar,bar width=16pt,
  symbolic x coords={parallel+,parallel-,short+,short-,steepener,flattener},xtick=data,
  x tick label style={font=\footnotesize},
  ylabel={change in EVE (USD billion)},
  tick label style={font=\small},label style={font=\small},
  ymin=-9.5,ymax=11,grid=major,grid style={gray!25},
  nodes near coords,every node near coord/.append style={font=\footnotesize,/pgf/number format/fixed,/pgf/number format/precision=2},
  table/col sep=comma]
\addplot[fill=omDef!60,draw=omDef] table[x=scenario,y=delta]{figdata/rates-credit-risk/24-liquidity-and-funding-risk-bank-treasury/eve.csv};
\draw[omThm,dashed,thick] (axis cs:{parallel+},-2.4) -- (axis cs:{flattener},-2.4);
\end{axis}
\end{tikzpicture}

\omcaption{Change in the stylised bank's \omterm{def:rc:liquidity-and-funding-risk-bank-treasury:eve}{economic value of equity} under the six supervisory
shocks, at a 1\% rate. The dashed line is the outlier threshold, a loss of 15\% of Tier~1 (2.4); the
parallel-up and steepener scenarios cross it. Data: the chapter's tutorial.}\label{fig:rc:liquidity-and-funding-risk-bank-treasury:eve}
\end{omfigure}

\section{Deposits: the behavioural model}

\begin{definition}[Non-maturity deposit, deposit beta]\label{def:rc:liquidity-and-funding-risk-bank-treasury:nmd}
A \emph{non-maturity deposit}\index{non-maturity deposit} can be withdrawn at any time and pays a rate
the bank sets; a behavioural model splits it into a core part, expected to stay and slotted over
several years, and a volatile part. Its \emph{deposit beta}\index{deposit beta} is the share of a
change in market rates passed on to the deposit rate.
\end{definition}

Deposit models turn a legally overnight liability into a long synthetic one: they justify
investing deposits in long assets. If the beta is low, rising rates widen the margin and NII rises;
if depositors demand a high beta, or leave, the margin reverses.

\begin{example}[Two betas]\label{ex:rc:liquidity-and-funding-risk-bank-treasury:beta}
The stylised bank earns NII of 4.27 a year at a 1\% rate. With a \omterm{def:rc:liquidity-and-funding-risk-bank-treasury:nmd}{deposit beta} of 0.35, a 200 basis
point rise adds 0.49 a year; with a beta of 0.80, the same rise costs 1.07
(\cref{fig:rc:liquidity-and-funding-risk-bank-treasury:nii}). The fixed-rate securities earn the same
in both cases; the difference is what the depositors take.
\end{example}

\begin{omfigure}
\begin{tikzpicture}
\begin{axis}[width=11.5cm,height=5.2cm,
  xlabel={parallel rise in rates (bp)},ylabel={change in NII (USD billion a year)},
  tick label style={font=\small},label style={font=\small},
  xmin=0,xmax=400,ymin=-2.4,ymax=1.2,grid=major,grid style={gray!25},
  legend columns=2,legend style={font=\footnotesize,at={(0.5,-0.25)},anchor=north,draw=none,fill=none,/tikz/every even column/.append style={column sep=8pt}},
  table/col sep=comma]
\addplot[omMeth,very thick] table[x=bp,y=beta35]{figdata/rates-credit-risk/24-liquidity-and-funding-risk-bank-treasury/nii.csv};
\addplot[omThm,very thick,dashed] table[x=bp,y=beta80]{figdata/rates-credit-risk/24-liquidity-and-funding-risk-bank-treasury/nii.csv};
\legend{deposit beta 0.35,deposit beta 0.80}
\end{axis}
\end{tikzpicture}

\omcaption{One-year \omterm{def:rc:liquidity-and-funding-risk-bank-treasury:eve}{net interest income sensitivity} of the stylised bank to a parallel rise in rates,
for two \omterm{def:rc:liquidity-and-funding-risk-bank-treasury:nmd}{deposit betas}. The sign of the bank's earnings exposure depends on how depositors behave.
Data: the chapter's tutorial.}\label{fig:rc:liquidity-and-funding-risk-bank-treasury:nii}
\end{omfigure}

\section{The 2023 regional-bank failures}

\begin{definition}[Held-to-maturity portfolio]\label{def:rc:liquidity-and-funding-risk-bank-treasury:htm}
A \emph{held-to-maturity portfolio}\index{held-to-maturity portfolio} holds securities the bank intends
and is able to keep until maturity, carried at amortised cost: their market losses do not appear in
earnings or in regulatory capital. Selling some can force the whole portfolio to be
marked.
\end{definition}

\begin{definition}[Bank run]\label{def:rc:liquidity-and-funding-risk-bank-treasury:run}
A \emph{bank run}\index{bank run} is a withdrawal of deposits faster than the bank can meet from its
liquid assets, driven by depositors' fear that others will withdraw first: individually rational,
collectively self-fulfilling.
\end{definition}

\begin{example}[The stylised bank a year later]\label{ex:rc:liquidity-and-funding-risk-bank-treasury:run}
After the 350 basis point rise, the MBS book held to maturity has lost 19.2\% of its value, USD~17.2
billion, 108\% of the bank's equity; the Treasuries have lost 3.9. On the day, the bank can pay out its
cash and sell its Treasuries at market: 36.1, 20.9\% of its deposits
(\cref{fig:rc:liquidity-and-funding-risk-bank-treasury:liquidity}). Anything more requires pledging or
selling the MBS, which realises the loss.
\end{example}

\begin{omfigure}
\begin{tikzpicture}
\begin{axis}[width=11.5cm,height=5.2cm,xbar stacked,bar width=18pt,
  ytick=\empty,xlabel={USD billion},
  tick label style={font=\small},label style={font=\small},
  xmin=0,xmax=190,ymin=0.4,ymax=1.9,
  legend columns=3,legend style={font=\footnotesize,at={(0.5,-0.3)},anchor=north,draw=none,fill=none,/tikz/every even column/.append style={column sep=8pt}}]
\addplot[fill=omMeth!60,draw=omMeth] coordinates {(15,1)};
\addplot[fill=omDef!60,draw=omDef] coordinates {(21.1,1)};
\addplot[fill=omExo!30,draw=omExo] coordinates {(72.8,1)};
\draw[omThm,very thick] (axis cs:173,0.5) -- (axis cs:173,1.5) node[above,font=\footnotesize] {deposits 173};
\draw[omThm,dashed,thick] (axis cs:36.1,0.5) -- (axis cs:36.1,1.5) node[above,font=\footnotesize] {same day 36.1};
\legend{cash,Treasuries at market,MBS at market (not same day)}
\end{axis}
\end{tikzpicture}

\omcaption{The stylised bank's liquid resources a year after the rise, at market values, against its
deposits. Only cash and Treasuries can be turned into payments on the day of a run: 20.9\% of
deposits. Data: the chapter's tutorial.}\label{fig:rc:liquidity-and-funding-risk-bank-treasury:liquidity}
\end{omfigure}

\begin{dated}{2026-09}{March 2023}\label{dat:rc:liquidity-and-funding-risk-bank-treasury:svb}
According to the Federal Reserve's review of April 2023, Silicon Valley Bank lost over 40 billion
dollars of deposits on 9 March 2023 and expected over 100 billion more on 10 March, roughly 85\% of its
deposit base; 94\% of its deposits had been uninsured at the end of 2022, and its \omterm{def:rc:liquidity-and-funding-risk-bank-treasury:htm}{held-to-maturity
portfolio} had a duration of 6.2 years. On 12 March the Federal Reserve created the Bank Term Funding
Program, lending for up to a year against eligible securities valued at par, and Silicon Valley Bank and
Signature Bank were resolved with all depositors protected. At the end of 2022 Silicon Valley Bank had
reported over 15 billion dollars of unrealised losses on its held-to-maturity securities, 89\% of its
common equity tier~1 capital (US Government Accountability Office). First Republic Bank, with 229.1 billion
dollars of assets, was closed on 1 May 2023 and its deposits and most assets sold to JPMorgan Chase; the
FDIC estimated the cost to its insurance fund at about 13 billion dollars.
\end{dated}

\section{Tutorial: a balance sheet under shocks}\label{tut:rc:liquidity-and-funding-risk-bank-treasury:tutorial}

\begin{tutorial}
\textbf{Goal.} Compute the stylised bank's EVE and NII sensitivities under the supervisory shocks, its
LCR, and the outflow its liquid assets can meet. \textbf{End state:} the numbers of
\cref{ex:rc:liquidity-and-funding-risk-bank-treasury:bank,ex:rc:liquidity-and-funding-risk-bank-treasury:eve,ex:rc:liquidity-and-funding-risk-bank-treasury:beta,ex:rc:liquidity-and-funding-risk-bank-treasury:run}
and the four charts.
\begin{tutsteps}
\item \textbf{Balance sheet}: \texttt{positions()}, \texttt{deposits(beta)}, \texttt{bank(rate)}.
\item \textbf{Value and earnings}: \texttt{eve\_table()}, \texttt{nii\_table()}.
\item \textbf{Liquidity}: \texttt{lcr\_now()}, \texttt{liquid\_same\_day()}, \texttt{htm\_loss()}.
\item \textbf{Charts}: \texttt{fig\_rc\_alm.py}.
\end{tutsteps}
\textbf{What to change next.} Hedge the MBS with a pay-fixed swap of the same duration and rerun the
EVE shocks; halve the core share of deposits and see the EVE change sign.
\end{tutorial}

\section{Build: the ALM engine}\label{bld:rc:liquidity-and-funding-risk-bank-treasury:alm}

\begin{build}
\bfield{Purpose} The firm's banking-book risk: EVE and NII under supervisory and internal scenarios,
liquidity ratios, transfer prices, and run capacity; the balance-sheet side of the risk engine.
\bfield{Interface} \texttt{shock\_curve(kind)}, \texttt{SCENARIOS}; \texttt{Position},
\texttt{Deposits}, \texttt{BalanceSheet} with \texttt{pv}, \texttt{delta\_eve}, \texttt{nii};
\texttt{market\_value}; \texttt{lcr}; \texttt{ftp\_rate}; \texttt{run\_capacity};
\texttt{bond\_price\_change}, \texttt{duration}; \texttt{eve\_ratio}.
\bfield{Rules} Flat base curve with scenario shifts; annual cash flows; core deposits slotted in equal
annual amounts; LCR with the Level~2 cap and the Level~2A haircut; outlier test at 15\% of Tier~1.
\bfield{Acceptance tests} \texttt{code/firm/alm/tests/}: shock shapes; bond price change against
duration; a matched book has no EVE sensitivity; the Level~2 cap; the FTP identity; run capacity.
\bfield{Stretch} A full yield curve with basis and optionality in loans and deposits; NSFR factors;
dynamic balance sheets; intraday liquidity.
\end{build}

\begin{omsources}
\item Board of Governors of the Federal Reserve System, \emph{Review of the Federal Reserve's
Supervision and Regulation of Silicon Valley Bank}, April 2023.
\item Federal Reserve press release of 12 March 2023 on the Bank Term Funding Program.
\item Basel Committee on Banking Supervision, \emph{The Liquidity Coverage Ratio and liquidity risk
monitoring tools} (2013); \emph{The net stable funding ratio} (2014); \emph{Interest rate risk in the
banking book} (2016).
\end{omsources}

\section{Exercises}

\begin{exercise}[$\star$]\label{exo:rc:liquidity-and-funding-risk-bank-treasury:1}
A bank has 30 of Level~1 assets and 40 of Level~2A assets before haircuts. What is its HQLA?
\end{exercise}

\begin{exercise}[$\star$]\label{exo:rc:liquidity-and-funding-risk-bank-treasury:2}
Estimate the price fall of a portfolio with duration 6.2 when yields rise 350 basis points, and
say why the exact figure differs.
\end{exercise}

\begin{exercise}[$\star$]\label{exo:rc:liquidity-and-funding-risk-bank-treasury:3}
What is the short-rate shock at five years in US dollars?
\end{exercise}

\begin{exercise}[$\star\star$]\label{exo:rc:liquidity-and-funding-risk-bank-treasury:4}
Why does the stylised bank gain EVE when rates fall, and what does that say about its position?
\end{exercise}

\begin{exercise}[$\star\star$]\label{exo:rc:liquidity-and-funding-risk-bank-treasury:5}
Why is a high share of uninsured deposits a liquidity risk and not only a credit issue for depositors?
\end{exercise}

\begin{exercise}[$\star\star$]\label{exo:rc:liquidity-and-funding-risk-bank-treasury:6}
How does an FTP curve change the incentive of a business that lends long?
\end{exercise}

\begin{exercise}[$\star\star\star$]\label{exo:rc:liquidity-and-funding-risk-bank-treasury:7}
\emph{Coding.} Compute the one-year NII change for a 350 basis point rise with \omterm{def:rc:liquidity-and-funding-risk-bank-treasury:nmd}{deposit betas} of 0.35
and 0.80.
\end{exercise}

\begin{exercise}[$\star\star\star$]\label{exo:rc:liquidity-and-funding-risk-bank-treasury:8}
\emph{Find the flaw.} ``Our securities are held to maturity and guaranteed by government agencies;
rate rises cannot hurt us because we will never sell them.''
\end{exercise}

\section{Problem: March 2023}

\begin{problem}[{Weekend problem --- a run on a solvent-looking bank}]\label{pb:rc:liquidity-and-funding-risk-bank-treasury:1}
The stylised bank of this chapter is a year into the rate rise. Its HTM portfolio is carried at cost;
its uninsured depositors are connected and follow the same news.

\medskip\noindent\textbf{Part I --- The balance sheet.}
\begin{enumerate}
\item Give the unrealised loss on the MBS and its share of equity.
\item Give the loss on the Treasuries.
\item What would the bank's equity be if both were marked?
\item Why did the supervisory EVE test already flag the risk at a 1\% rate?
\item What hedge would have removed most of it, and at what cost to earnings?
\end{enumerate}

\medskip\noindent\textbf{Part II --- Liquidity.}
\begin{enumerate}[resume]
\item Give the LCR and why it is below 100\%.
\item Give the outflow the bank can meet on the day.
\item Compare it with the first day of SVB's run as a share of deposits.
\item What sources of funding could the bank use on day two?
\item How did the Bank Term Funding Program change the bank's options?
\end{enumerate}

\medskip\noindent\textbf{Part III --- Behaviour.}
\begin{enumerate}[resume]
\item Why is a 94\% uninsured deposit base different from a retail base of the same size?
\item How does a \omterm{def:rc:liquidity-and-funding-risk-bank-treasury:nmd}{deposit beta} rising with rates change NII and EVE?
\item What does announcing a securities sale and a capital raise signal to depositors?
\item Why can a run happen to a bank that is solvent at amortised cost?
\item Which ratio would have measured the concentration of depositors?
\end{enumerate}

\medskip\noindent\textbf{Part IV --- Judgement.}
\begin{enumerate}[resume]
\item What should the treasury have done in 2021?
\item Should held-to-maturity losses count in capital?
\item How should deposit models be validated?
\item State the \emph{named result}: the held-to-maturity unrealised loss as a share of capital and
the one-day outflow at which the \omterm{def:rc:liquidity-and-funding-risk-bank-treasury:run}{bank runs} out of liquid assets.
\item In one sentence: what killed the bank, market risk or liquidity risk?
\end{enumerate}
\end{problem}

\section{Interview questions}

\begin{interviewq}[$\star$ \iqroles{bank, risk}]\label{iq:rc:liquidity-and-funding-risk-bank-treasury:1}
Define the LCR and the NSFR. What does each protect against?
\end{interviewq}

\begin{interviewq}[$\star\star$ \iqroles{bank}]\label{iq:rc:liquidity-and-funding-risk-bank-treasury:2}
What is \omterm{def:rc:liquidity-and-funding-risk-bank-treasury:ftp}{funds-transfer pricing} and why do banks need it?
\end{interviewq}

\begin{interviewq}[$\star\star$ \iqroles{risk, researcher}]\label{iq:rc:liquidity-and-funding-risk-bank-treasury:3}
EVE and NII sensitivity can have opposite signs. Explain and give an example.
\end{interviewq}

\begin{interviewq}[$\star\star$ \iqroles{researcher}]\label{iq:rc:liquidity-and-funding-risk-bank-treasury:4}
How would you model \omterm{def:rc:liquidity-and-funding-risk-bank-treasury:nmd}{non-maturity deposits}?
\end{interviewq}

\begin{interviewq}[$\star\star\star$ \iqroles{bank, risk}]\label{iq:rc:liquidity-and-funding-risk-bank-treasury:5}
What went wrong at Silicon Valley Bank, in risk-management terms?
\end{interviewq}

\begin{interviewq}[$\star\star\star$ \iqroles{trader, bank}]\label{iq:rc:liquidity-and-funding-risk-bank-treasury:6}
How would you hedge a bank's EVE exposure, and what does the hedge do to earnings and accounting?
\end{interviewq}
